% !TEX program = pdflatex
% !TEX root = p_ALLP.tex
\documentclass[12pt]{article}
\usepackage{amsfonts}
\usepackage{mathrsfs}
\usepackage{amsthm,amssymb,amscd,amsmath}
\usepackage{tikz-cd}
\textheight=21.4 cm \textwidth=15cm  \topmargin=-0.2 in
\oddsidemargin 0cm \evensidemargin -0.2cm
\parindent 24pt
\baselineskip 4pt

\newtheorem{theorem}{Theorem}[section]
\newtheorem{lemma}[theorem]{Lemma}
\newtheorem{corollary}[theorem]{Corollary}
\newtheorem{definition}[theorem]{Definition}
\newtheorem{proposition}[theorem]{Proposition}
\newtheorem{remark}[theorem]{Remark}

\newcommand{\pcb}{pcb}
\newcommand{\CBp}{CB_p}
\newcommand{\otimesp}{\mathop{\overset{\wedge p}{\otimes}}}

\title{Local Lifting Properties and Injectivity for $p$-Operator Spaces}
\author{Yafei Zhao and Zhe Dong\thanks{Corresponding author.}}
\date{}

\begin{document}
\maketitle

\begin{center}
\begin{minipage}{11.5cm}
{\bf Abstract:} For $1<p<\infty$, we investigate local lifting properties and injectivity for $p$-operator spaces. We prove that for any $p$-operator space, the $\lambda$-$p$-approximating local lifting property is equivalent to the $\lambda$-$p$-local lifting property. For dual $p$-operator spaces, we show that approximate $\lambda$-$p$-injectivity coincides with $\lambda$-$p$-injectivity. In particular, we show that every $n$-dimensional $p$-operator space is $n$-$p$-injective and has the $n$-$p$-LLP.
\par
{\bf Keywords:} $p$-operator space, $\lambda$-$p$-approximating local lifting property,
approximate $\lambda$-$p$-injectivity.
\par
{\bf 2020 MR Subject Classification:} Primary 46L07.
\end{minipage}
\end{center}

\section{Introduction}

The local theory of operator spaces is a non-commutative quantization of the local theory of Banach spaces initiated by Grothendieck \cite{Gr55}. In the classical setting of operator spaces, the abstract matrix norm characterization was established by Ruan \cite{Ru89}, and the corresponding dual matrix norm theory was developed by Blecher \cite{Ble92}. Local properties such as local reflexivity, exactness, nuclearity, injectivity and the local lifting property play a central role in operator space theory. Building on Connes' fundamental work \cite{Co76}, Choi and Effros \cite{CE76,CE77b} established that a $C^*$-algebra $\mathcal{A}$ is nuclear if and only if its second dual $\mathcal{A}^{**}$ is injective. Effros, Ozawa, and Ruan \cite{EOR01} generalized this to arbitrary operator spaces, proving that an operator space $V$ is nuclear if and only if $V$ is locally reflexive and $V^{**}$ is injective. Dong and Ruan \cite{DR07} showed that an operator space $V$ is exact if and only if $V$ is locally reflexive and $V^{**}$ is weak$^*$ exact.

In this paper, we study the local lifting properties and injectivity for $p$-operator spaces. Let us recall from the classical operator space theory \cite{KyeRuan1999} that an operator space $V$ is said to have the $\lambda$-local lifting property ($\lambda$-LLP) if given any operator spaces $W \subseteq Y$, a complete contraction $\varphi: V \to Y/W$, any finite-dimensional subspace $E \subseteq V$, and $\varepsilon>0$, there exists a completely bounded map $\tilde{\varphi}: E \to Y$ such that
\[
q\circ \tilde{\varphi}=\varphi|_E
\qquad\text{and}\qquad
\|\tilde{\varphi}\|_{\mathrm{cb}}< \lambda+\varepsilon,
\]
where $q: Y \to Y/W$ is the complete quotient map. If $\lambda = 1$, we simply say that $V$ has the LLP. An operator space $V$ is \emph{$\lambda$-injective} if given operator spaces $W\subseteq Y$, and any completely bounded map $\psi:W\to V$, there is an extension $\widetilde{\psi}:Y\to V$ such that
\[
\widetilde{\psi}|_{W}=\psi,\qquad
\|\widetilde{\psi}\|_{\mathrm{cb}}\leq\lambda\|\psi\|_{\mathrm{cb}}.
\]
For $\lambda=1$, we say that $V$ is injective. The local lifting property is a natural non-commutative analogue of the Banach space local lifting property, but it exhibits much richer structures. Kye and Ruan \cite{KyeRuan1999} proved that an operator space $V$ has the $\lambda$-LLP if and only if its operator dual $V^*$ is $\lambda$-injective. This duality theorem reveals a deep connection between the local lifting property and injectivity.

Let $1<p<\infty$. Following Pisier \cite{Pisier1990} and Le Merdy \cite{LeMerdy1996}, Daws \cite{Daws2007} presented the abstract framework for $p$-operator spaces, which generalizes classical operator spaces (the $p=2$ case) to the $L^p$ setting. The corresponding theory of $p$-operator spaces has since been developed in considerable detail. The $p$-approximation property was investigated by An et al.\ \cite{AnLeeRuan2010}. Lee \cite{Lee2014} studied the local structure of $p$-operator spaces through conditions $C_p$, $C'_p$, and $C''_p$. Zhao and Dong \cite{ZD14} introduced $p$-mapping spaces and applied them to the study of $p$-local reflexivity. Moreover, if $V$ is a $p$-operator space, then its dual space $V^* = \CBp(V, \mathbb{C})$ is naturally a $p$-operator space and agrees isometrically with the Banach dual (see \cite{Daws2007}).

In the study of local properties, it is natural to consider the approximate versions of these properties. For classical operator spaces, Dong \cite{Dong2007} introduced the approximating local lifting property (ALLP) and approximate injectivity, and showed that for an operator space $V$,
\[
\begin{aligned}
V\text{ has the ALLP} &\Longleftrightarrow V\text{ has the LLP} \\
&\Longleftrightarrow V^*\text{ is approximately injective} \\
&\Longleftrightarrow V^*\text{ is injective}.
\end{aligned}
\]
Dong \cite{Dong2009} further investigated local lifting properties of operator spaces. This motivates us to investigate whether similar equivalence holds in the category of $p$-operator spaces. However, the classical proof of the equivalence ($V\text{ has the }\lambda\text{-}p\text{-LLP} \Longleftrightarrow V^*\text{ is }\lambda\text{-}p\text{-injective}$) does not extend to $p$-operator spaces, since the Arveson--Wittstock extension theorem fails in general \cite{Lee2015}. The main purpose of this paper is to answer whether the $\lambda$-$p$-ALLP implies the $\lambda$-$p$-LLP, and whether approximate $\lambda$-$p$-injectivity implies $\lambda$-$p$-injectivity for dual $p$-operator spaces.

In Section 2, we establish the equivalence between the $\lambda$-$p$-ALLP and the $\lambda$-$p$-LLP for arbitrary $p$-operator spaces (Theorem~\ref{thm:main}). The proof relies on quantitative estimates for finite-dimensional $p$-operator spaces. In Section 3, we prove that for dual $p$-operator spaces, approximate $\lambda$-$p$-injectivity implies $\lambda$-$p$-injectivity (Theorem~\ref{thm:dual-exact-approx}).  We also use a complement projection trick for extensions into $V^*$ (Proposition~\ref{prop:complement-extension}) and record the resulting two-step scheme (Corollary~\ref{cor:llp-vstar-two-step} and Remark~\ref{rem:llp-to-approx-inj-open}).  When $p=2$, for each $\lambda\geq 1$,
\[
\begin{aligned}
V\text{ has the $\lambda$-ALLP} &\Longleftrightarrow V\text{ has the $\lambda$-LLP} \\
&\Longleftrightarrow V^*\text{ is approximately $\lambda$-injective} \\
&\Longleftrightarrow V^*\text{ is $\lambda$-injective}
\end{aligned}
\]
(Corollary~\ref{cor:dual-p-2}). Finally, in Section 4, we apply these results to finite-dimensional $p$-operator spaces. We show that every $n$-dimensional $p$-operator space is $n$-$p$-injective and has the $n$-$p$-LLP (Corollary~\ref{cor:fd-equivalence}).

\section{$\lambda$-$p$-Local lifting properties}

\begin{definition}\label{def:lambda-p-lp}
Let $\lambda\geq 1$.  A $p$-operator space $V$ has the
\emph{$\lambda$-$p$-lifting property} ($\lambda$-$p$-LP) if given $p$-operator spaces $W\subseteq Y$, a $p$-complete contraction
$\varphi:V\to Y/W$, and $\varepsilon>0$, there exists a $p$-completely bounded
linear map $\widetilde{\varphi}:V\to Y$ such that
\[
        q\circ\widetilde{\varphi}=\varphi
        \qquad\text{and}\qquad
        \|\widetilde{\varphi}\|_{\pcb}<\lambda+\varepsilon,
\]
i.e., we have the commutative diagram
\begin{equation*}
\begin{tikzcd}
& Y \arrow{d}{q} \\
V \arrow{r}[swap]{\varphi} \arrow{ru}{\widetilde{\varphi}} & Y/W\text{,}
\end{tikzcd}
\end{equation*}
where $q:Y\to Y/W$ denotes the
$p$-complete quotient map from $Y$ onto $Y/W$.
For $\lambda=1$, we say that $V$ has the $p$-LP.
When $p=2$, we say that $V$ has the $\lambda$-LP.
\end{definition}

\begin{definition}\label{def:lambda-llp}
Let $\lambda\geq 1$.  A $p$-operator space $V$ has the
\emph{$\lambda$-$p$-local lifting property} ($\lambda$-$p$-LLP) if given $p$-operator spaces $W\subseteq Y$, a $p$-complete contraction
$\varphi:V\to Y/W$, for every finite-dimensional subspace $E\subseteq V$, and $\varepsilon>0$, there exists a $p$-completely bounded
linear map $\widetilde{\varphi}:E\to Y$ such that
\[
        q\widetilde{\varphi}=\varphi|_E
        \qquad\text{and}\qquad
        \|\widetilde{\varphi}\|_{\pcb}<\lambda+\varepsilon,
\]
i.e., we have the commutative diagram
\begin{equation*}
\begin{tikzcd}
& & Y \arrow{d}{q} \\
E \arrow[hook]{r} \arrow{rru}{\widetilde{\varphi}} & V \arrow{r}{\varphi} & Y/W\text{,}
\end{tikzcd}
\end{equation*}
where $q:Y\to Y/W$ denotes the
$p$-complete quotient map from $Y$ onto $Y/W$.
For $\lambda=1$, we say that $V$ has the $p$-LLP.
When $p=2$, we say that $V$ has the $\lambda$-LLP.
\end{definition}

\begin{definition}\label{def:lambda-allp}
Let $\lambda\geq 1$.  A $p$-operator space $V$ has the
\emph{$\lambda$-$p$-approximating local lifting property} ($\lambda$-$p$-ALLP) if
given $p$-operator spaces $W\subseteq Y$, and a $p$-complete contraction
$\varphi:V\to Y/W$, for every finite-dimensional subspace $E\subseteq V$, and
$\varepsilon>0$, there exists a net $\{\varphi_\alpha\}$ of $p$-completely bounded
linear maps $\varphi_\alpha:E\to Y$ such that
$\|\varphi_\alpha\|_{\pcb}<\lambda+\varepsilon$
and $q\varphi_\alpha\to \varphi|_E$ in the point-norm topology on $Y/W$, i.e., we
have the approximately commutative diagram
\begin{equation*}
\begin{tikzcd}
& & Y \arrow{d}{q} \\
E \arrow[hook]{r} \arrow{rru}{\varphi_\alpha} & V \arrow{r}{\varphi} & Y/W\text{,}
\end{tikzcd}
\end{equation*}
where $q:Y\to Y/W$ denotes the
$p$-complete quotient map from $Y$ onto $Y/W$.
For $\lambda=1$, we say that $V$ has the $p$-ALLP.
When $p=2$, we say that $V$ has the $\lambda$-ALLP.
\end{definition}

From the definitions, for every $p$-operator space $V$ and every $\lambda\geq 1$ we have
\[
        \lambda\hbox{-}p\hbox{-LP}
        \quad\Longrightarrow\quad
        \lambda\hbox{-}p\hbox{-LLP}
        \quad\Longrightarrow\quad
        \lambda\hbox{-}p\hbox{-ALLP}.
\]
To show that the $\lambda$-$p$-LLP and the $\lambda$-$p$-ALLP are equivalent, we first prove two quantitative estimates for finite-dimensional $p$-operator spaces.

\begin{lemma}\label{lem:finite-cb-comparison}
Let $V$ and $W$ be $p$-operator spaces, and suppose that either $V$ or $W$ is
$n$-dimensional.  Then every linear map $\varphi:V\to W$ is $p$-completely
bounded and
\[
        \|\varphi\|_{\pcb}\leq n\|\varphi\|.
\]
\end{lemma}

\begin{proof}
The proof adapts the argument for classical operator spaces (see \cite[Corollary 2.2.4]{ER00}).
Let us suppose that $W$ is $n$-dimensional. Let $w_1,\ldots,w_n$ be an Auerbach
basis for $W$, with corresponding coordinate functionals $g_1,\ldots,g_n \in W^*$, so that
$\|w_k\|=\|g_k\|=1$. Since
\[
        \mathrm{id}_W = \sum_{k=1}^n \theta_{w_k} \circ g_k,
\]
where $\theta_{w_k}(\alpha) = \alpha w_k$ are $p$-complete isometries from $\mathbb{C}$ into $W$, we have
\[
        \varphi = \sum_{k=1}^n \theta_{w_k} \circ g_k \circ \varphi.
\]
By \cite[Lemma 4.2]{Daws2007}, we have that $g_k \circ \varphi$ is a $p$-completely bounded linear functional with
\[
        \|g_k \circ \varphi\|_{\pcb} = \|g_k \circ \varphi\| \leq \|\varphi\|.
\]
It follows that
\[
        \|\varphi\|_{\pcb} \leq \sum_{k=1}^n \|\theta_{w_k}\|_{\pcb}\|g_k \circ \varphi\|_{\pcb} = \sum_{k=1}^n \|g_k \circ \varphi\| \leq n\|\varphi\|.
\]

The case where $V$ is $n$-dimensional follows similarly by applying the same argument to $\varphi:V\to\varphi(V)$.
\end{proof}

\begin{lemma}\label{lem:correction}
Let $q:Y\to Z$ be a $p$-complete quotient map of $p$-operator spaces, and let
$E$ be a finite-dimensional $p$-operator space with $n=\dim E\geq 1$.  Then
the induced map $\rho:\CBp(E,Y)\longrightarrow \CBp(E,Z)$ is a surjective
$p$-complete contraction.
Moreover, for each $\phi\in \CBp(E,Z)$ and $\varepsilon>0$, there exists $\psi\in\CBp(E,Y)$ such that
\[
        q\psi=\phi,\qquad \|\psi\|_{\pcb}\leq n(1+\varepsilon)\|\phi\|_{\pcb}.
\]
\end{lemma}

\begin{proof}
Post-composition with a $p$-complete quotient map is a $p$-complete contraction.

We let $\phi\in\CBp(E,Z)$. We can choose an Auerbach basis $e_1,\ldots,e_n$ of $E$ with coordinate functionals $f_1,\ldots,f_n\in E^*$, so that $\|e_k\|=\|f_k\|=1$. Then $\|\phi(e_k)\|\leq\|\phi\|\leq\|\phi\|_{\pcb}$. Since $q$ is a $p$-complete quotient map, for any $\varepsilon>0$ we can choose $y_k\in Y$ such that $q(y_k)=\phi(e_k)$ and
\[
        \|y_k\|\leq (1+\varepsilon)\|\phi(e_k)\|\leq (1+\varepsilon)\|\phi\|_{\pcb}.
\]
We define $\psi:E\to Y$ by $\psi(e_k)=y_k$. For each $k$, let $\theta_{y_k}:\mathbb{C}\to Y$ be given by $\theta_{y_k}(\alpha)=\alpha y_k$. Then $\psi=\sum_{k=1}^n \theta_{y_k}\circ f_k$, and $\|\theta_{y_k}\|_{\pcb}=\|y_k\|$ for each $k$.
Then $q\circ\psi=\phi$, which shows that post-composition with $q$ is surjective.

Moreover, each $f_k$ is a linear functional, so $\|f_k\|_{\pcb}=\|f_k\|=1$ by \cite[Lemma 4.2]{Daws2007}.
Hence
\[
        \|\psi\|_{\pcb}
        \leq \sum_{k=1}^n \|\theta_{y_k}\|_{\pcb}\,\|f_k\|_{\pcb}
        =\sum_{k=1}^n \|y_k\|
        \leq n(1+\varepsilon)\|\phi\|_{\pcb}.
\]
\end{proof}

\begin{theorem}\label{thm:main}
Let $V$ be a $p$-operator space and let $\lambda\geq 1$.  Then $V$ has the
$\lambda$-$p$-ALLP if and only if $V$ has the $\lambda$-$p$-LLP.
\end{theorem}

\begin{proof}
It is clear that the $\lambda$-$p$-LLP implies the $\lambda$-$p$-ALLP.

Conversely, we assume that $V$ has the $\lambda$-$p$-ALLP.  We suppose that $W\subseteq Y$ are $p$-operator spaces. We denote by $q:Y\to Y/W$ the $p$-complete quotient map.  Let $\varphi:V\to Y/W$ be a $p$-complete contraction, let $E\subseteq V$ be finite-dimensional, and let $\varepsilon>0$.  If $\varphi|_E=0$, then the zero operator $0:E\to Y$ satisfies $q\circ 0=\varphi|_E$. This operator is the required lifting.  We may therefore assume that $\varphi|_E\neq 0$.

We let $n=\dim E$.  From Lemma~\ref{lem:finite-cb-comparison}, we have
$\|\sigma\|_{\pcb}\leq n\|\sigma\|$ for any linear map $\sigma:E\to Y/W$.
We choose $\delta>0$ small enough that
\[
        n^{2}\delta<\frac{\varepsilon}{2}.
\]
By the $\lambda$-$p$-ALLP applied to $\varphi$ with $\varepsilon/2$, there
is a net $\{w_\alpha\}$ in $\CBp(E,Y)$ with
\[
        \|w_\alpha\|_{\pcb}<\lambda+\frac{\varepsilon}{2}.
\]
and $qw_\alpha\to \varphi|_E$ in the point-norm topology on $Y/W$.  Since $E$ is
finite-dimensional, point-norm convergence is equivalent to operator norm
convergence on $E$, so we may choose $w_0=w_\alpha$ satisfying
\[
        \|qw_0-\varphi|_E\|<\delta.
\]
We define $\varphi_0=w_0$.  Then we have
\[
        \|\varphi_0\|_{\pcb}<\lambda+\frac{\varepsilon}{2}
\]
and
\[
        \|q\varphi_0-\varphi|_E\|<\delta.
\]
From Lemma~\ref{lem:finite-cb-comparison} we thus have
$\|q\varphi_0-\varphi|_E\|_{\pcb}\leq n\delta$.  We define
\[
        \phi=\varphi|_E-q\varphi_0:E\to Y/W.
\]
From Lemma~\ref{lem:correction}, for any $\varepsilon'>0$ there exists $\psi:E\to Y$ with $q\psi=\phi$ and $\|\psi\|_{\pcb}\leq n(1+\varepsilon')\|\phi\|_{\pcb}$. Since $n^{2}\delta<\frac{\varepsilon}{2}$, we can choose $\varepsilon'$ small enough so that
\[
        \|\psi\|_{\pcb}\leq n(1+\varepsilon')\|\phi\|_{\pcb}\leq n(1+\varepsilon')n\delta <\frac{\varepsilon}{2}.
\]
We define
\[
        \widetilde{\varphi}=\varphi_0+\psi:E\to Y.
\]
Then
\[
        q\widetilde{\varphi}=q\varphi_0+q\psi=q\varphi_0+(\varphi|_E-q\varphi_0)=\varphi|_E,
\]
and
\[
        \|\widetilde{\varphi}\|_{\pcb}
        \leq \|\varphi_0\|_{\pcb}+\|\psi\|_{\pcb}
        <\lambda+\varepsilon.
\]
This shows that $V$ has the $\lambda$-$p$-LLP.
\end{proof}

\begin{remark}\label{rem:main-on-Lp}
The proof of Theorem~\ref{thm:main} uses only
Definitions~\ref{def:lambda-llp}--\ref{def:lambda-allp} and
Lemma~\ref{lem:finite-cb-comparison}.
Hence the same argument shows that for every $p$-operator space $V$ and every
$\lambda\geq 1$, the $\lambda$-$p$-ALLP and the $\lambda$-$p$-LLP are equivalent
whenever in those definitions one restricts to $p$-operator spaces
$W\subseteq Y$ on $L_p$ spaces.
\end{remark}

\begin{proposition}\label{prop:complemented}
Let $V$ be a $p$-operator space and let $P:V\to V$ be a $p$-completely bounded
projection.  We denote by $V_0$ the subspace $P(V)\subseteq V$, with the
$p$-operator space structure induced by $V$.
\begin{enumerate}
        \item[\rm(i)] If $V$ has the $\lambda$-$p$-LLP, then $V_0$ has the
                $\lambda\|P\|_{\pcb}$-$p$-LLP.
        \item[\rm(ii)] If $V$ has the $\lambda$-$p$-ALLP, then $V_0$ has the
                $\lambda\|P\|_{\pcb}$-$p$-ALLP.
\end{enumerate}
In particular, whenever $P$ is a $p$-complete contraction, $V_0$ has the
$\lambda$-$p$-LLP in case {\rm(i)} and the $\lambda$-$p$-ALLP in case {\rm(ii)}.
\end{proposition}

\begin{proof}
Part~{\rm(ii)} follows from part~{\rm(i)} and Theorem~\ref{thm:main}.  We prove
{\rm(i)}.  We suppose that $W\subseteq Y$ are $p$-operator spaces. We denote by $q:Y\to Y/W$ the $p$-complete quotient map.  Let $\varphi:V_0\to Y/W$ be a $p$-complete contraction and let $E\subseteq V_0$ be finite-dimensional.  If $\varphi P=0$, then $\varphi|_E=0$ and the zero operator $0:E\to Y$ satisfies $q\circ 0=\varphi|_E$. This operator is a lifting of $\varphi|_E$.  If $\varphi P\neq 0$, we normalize the composition
\[
        \theta=\frac{\varphi P}{\|\varphi P\|_{\pcb}}:V\longrightarrow Y/W.
\]
Since $P|_{V_0}=\mathrm{id}_{V_0}$, we have $(\varphi P)|_E=\varphi|_E$, and
\[
        \|\varphi P\|_{\pcb}\leq \|\varphi\|_{\pcb}\|P\|_{\pcb}.
\]
Applying the $\lambda$-$p$-LLP of $V$ to the finite-dimensional subspace
$E\subseteq V$ and the $p$-complete contraction $\theta$, for any $\delta>0$ we obtain
$\widetilde{\theta}:E\to Y$ such that
\[
        q\widetilde{\theta}=\theta|_E
\]
and
\[
        \|\widetilde{\theta}\|_{\pcb}<\lambda+\delta.
\]
We define $\widetilde{\varphi}=\|\varphi P\|_{\pcb}\widetilde{\theta}$.  Then
$q\widetilde{\varphi}=\varphi|_E$ and
\[
        \|\widetilde{\varphi}\|_{\pcb}
        <(\lambda+\delta)\|\varphi P\|_{\pcb}
        \leq(\lambda+\delta)\|P\|_{\pcb}.
\]
Given $\varepsilon>0$, we choose $\delta>0$ so that
$\delta\|P\|_{\pcb}<\varepsilon$.  This is the desired estimate for
the constant $\lambda\|P\|_{\pcb}$.
\end{proof}

\section{$\lambda$-$p$-injectivity and duality}

\begin{definition}\label{def:p-injective}
Let $\lambda\geq 1$.  A $p$-operator space $V$ is
\emph{$\lambda$-$p$-injective} if given $p$-operator spaces $W_0\subseteq W$, and any $p$-completely bounded map $\varphi:W_0\to V$, there is an extension
$\widetilde{\varphi}:W\to V$ such that
\[
        \widetilde{\varphi}|_{W_0}=\varphi,\qquad
        \|\widetilde{\varphi}\|_{\pcb}
        \leq\lambda\|\varphi\|_{\pcb}.
\]
For $\lambda=1$, we say that $V$ is $p$-injective.
When $p=2$, we say that $V$ is $\lambda$-injective.
\end{definition}

\begin{definition}\label{def:approx-p-injective}
Let $\lambda\geq 1$.  A $p$-operator space $V$ is
\emph{approximately $\lambda$-$p$-injective} if given $p$-operator spaces $W_0\subseteq W$, and any $p$-completely bounded map $\varphi:W_0\to V$, for any $\varepsilon>0$
there exists a net $\{\widetilde{\varphi}_\alpha\}$ of $p$-completely bounded
linear maps $\widetilde{\varphi}_\alpha:W\to V$ such that
\[
        \|\widetilde{\varphi}_\alpha\|_{\pcb}
        \leq(\lambda+\varepsilon)\|\varphi\|_{\pcb}
\]
and $\widetilde{\varphi}_\alpha|_{W_0}\to\varphi$ in the point-norm topology on
$W_0$.
For $\lambda=1$, we say that $V$ is approximately $p$-injective.
When $p=2$, we say that $V$ is approximately $\lambda$-injective.
\end{definition}

\begin{remark}
For classical operator spaces ($p=2$), Dong's original definition of approximate injectivity \cite{Dong2007} requires the approximating net to be completely contractive. The $\varepsilon$-slack in Definition \ref{def:approx-p-injective} serves as a natural quantitative relaxation that parallels the definition of the $\lambda$-$p$-ALLP.
\end{remark}

It is immediate from Definitions~\ref{def:p-injective} and~\ref{def:approx-p-injective} that $\lambda$-$p$-injectivity implies approximate $\lambda$-$p$-injectivity.

The scalar-valued extension theorem remains available in the $p$-operator
space setting, although the Arveson--Wittstock extension theorem does not
extend to arbitrary ranges.

As noted in the introduction, the classical equivalence between the local lifting property of $V$ and the injectivity of $V^*$ relies on the Arveson--Wittstock extension theorem. The scalar-valued extension theorem alone is much weaker; it survives in the $p$-operator space setting, but it does not by itself recover the full duality.

\begin{proposition}\label{prop:equivalence-extension}
Let $1 < p < \infty$. Suppose that, for every $p$-operator space $V$,
$V$ has the $p$-LLP if and only if $V^*$ is $p$-injective.  Then the
$p$-completely bounded scalar extension theorem holds: whenever
$W_0\subseteq W$ are $p$-operator spaces, every $p$-complete contraction
$\psi:W_0\to\mathbb{C}$ admits a $p$-completely contractive extension
$\widetilde{\psi}:W\to\mathbb{C}$.
\end{proposition}

\begin{proof}
First, $\mathbb{C}$ has the $p$-LLP.  Indeed, let $Y_0 \subseteq Y$ be
$p$-operator spaces, let $q:Y\to Y/Y_0$ be the $p$-complete quotient map, and
let $\varphi:\mathbb{C}\to Y/Y_0$ be a $p$-complete contraction.  Given
$\varepsilon>0$, choose $y\in Y$ such that $q(y)=\varphi(1)$ and
$\|y\|<1+\varepsilon$.  The map
$\widetilde{\varphi}(\alpha)=\alpha y$ satisfies
$q\widetilde{\varphi}=\varphi$ and
$\|\widetilde{\varphi}\|_{\pcb}=\|y\|<1+\varepsilon$.

By the assumed duality, $\mathbb{C}^*\cong\mathbb{C}$ is $p$-injective.  Hence
each $p$-complete contraction $\psi:W_0\to\mathbb{C}$ extends to a
$p$-completely bounded map $\widetilde{\psi}:W\to\mathbb{C}$ with
$\|\widetilde{\psi}\|_{\pcb}\leq 1$.  This is the scalar extension theorem.
\end{proof}

\begin{proposition}\label{prop:scalar-injective}
Let $1<p<\infty$.  Then $\mathbb{C}$ is $p$-injective.
\end{proposition}

\begin{proof}
Let $W_0\subseteq W$ be $p$-operator spaces and let $\psi:W_0\to\mathbb C$
be a $p$-completely bounded linear functional.  By
\cite[Lemma~4.2]{Daws2007}, every bounded linear functional on a $p$-operator
space is $p$-completely bounded with $p$-completely bounded norm equal to its
ordinary norm; applied to $\psi$, this gives
$\|\psi\|_{\pcb}=\|\psi\|$.  The Hahn--Banach theorem provides a norm-preserving
linear extension $\widetilde\psi:W\to\mathbb C$ with
$\|\widetilde\psi\|=\|\psi\|$, and the same lemma applied to $\widetilde\psi$
yields $\|\widetilde\psi\|_{\pcb}=\|\widetilde\psi\|$.  Hence
\[
        \|\widetilde\psi\|_{\pcb}
        =\|\widetilde\psi\|
        =\|\psi\|
        =\|\psi\|_{\pcb},
\]
so $\mathbb{C}$ is $p$-injective.
\end{proof}

\begin{proposition}\label{prop:complement-extension}
Let $V$ be a $p$-operator space, let $W_0\subseteq W$ be $p$-operator spaces,
and let $P:W\to W$ be a $p$-completely bounded projection with $P(W)=W_0$ and
$P|_{W_0}=\mathrm{id}_{W_0}$.  If $\psi:W_0\to V^*$ is $p$-completely bounded,
then $\Psi:=\psi\circ P:W\to V^*$ satisfies $\Psi|_{W_0}=\psi$ and
\[
        \|\Psi\|_{\pcb}\leq \|P\|_{\pcb}\,\|\psi\|_{\pcb}.
\]
In particular, if $\|P\|_{\pcb}\leq 1$, then every $p$-completely bounded map
$W_0\to V^*$ extends to $W$ without enlarging the $p$-completely bounded norm.
\end{proposition}

\begin{proof}
For $w_0\in W_0$ we have $\Psi(w_0)=\psi(Pw_0)=\psi(w_0)$.  The norm estimate
follows from $\|\Psi\|_{\pcb}=\|\psi\circ P\|_{\pcb}
\leq\|\psi\|_{\pcb}\|P\|_{\pcb}$.
\end{proof}

\begin{corollary}\label{cor:llp-vstar-two-step}
Let $V$ be a $p$-operator space and let $\lambda\geq 1$.  Suppose that $V$
has the $\lambda$-$p$-LLP and that $V^*$ is approximately $\lambda$-$p$-injective.
Then $V^*$ is $\lambda$-$p$-injective.
\end{corollary}

\begin{proof}
This is immediate from Theorem~\ref{thm:dual-exact-approx}.
\end{proof}

\begin{remark}\label{rem:llp-to-approx-inj-open}
Corollary~\ref{cor:llp-vstar-two-step} isolates the part of the classical
duality chain that still works in the $p$-operator space setting once one has
approximate information on $V^*$: Theorem~\ref{thm:dual-exact-approx} upgrades
approximate $\lambda$-$p$-injectivity of $V^*$ to $\lambda$-$p$-injectivity.
The remaining step,
\[
        V\text{ has the }\lambda\text{-}p\text{-LLP}
        \quad\Longrightarrow\quad
        V^*\text{ is approximately }\lambda\text{-}p\text{-injective},
\]
does not follow from the quotient-side dual identification in \cite{Daws2007}
alone.  When $p=2$,
it is part of the known equivalences in \cite{KyeRuan1999,Dong2007} and is
recorded in Corollary~\ref{cor:dual-p-2}.  For general $p$, we leave this
implication as an open problem.  Proposition~\ref{prop:complement-extension}
gives a concrete situation in which extensions into $V^*$ are available without
any quotient identification of duals.
\end{remark}

\begin{proposition}\label{prop:proj-injective}
Let $V$ be a $p$-operator space and let $P:V\to V$ be a $p$-completely bounded
projection.  We denote by $V_0$ the subspace $P(V)\subseteq V$, with the
$p$-operator space structure induced by $V$.
\begin{enumerate}
        \item[\rm(i)] If $V$ is $\lambda$-$p$-injective, then $V_0$ is
                $\lambda\|P\|_{\pcb}$-$p$-injective.
        \item[\rm(ii)] If $V$ is approximately $\lambda$-$p$-injective, then $V_0$
                is approximately $(\lambda\|P\|_{\pcb})$-$p$-injective.
\end{enumerate}
In particular, if $P$ is a $p$-complete contraction, then $V_0$ is
(approximately) $\lambda$-$p$-injective.
\end{proposition}

\begin{proof}
We prove (ii). Part (i) is obtained by specializing to $\varepsilon=0$.

Given $p$-operator spaces $W_0\subseteq W$, and a $p$-completely bounded map $\varphi:W_0\to V_0$, we regard $\varphi$ as a
$p$-completely bounded map $\varphi:W_0\to V$.  By approximate $\lambda$-$p$-injectivity of $V$,
for any $\eta>0$ there is a net $\{\widetilde{\varphi}_\alpha\}$ of maps
$\widetilde{\varphi}_\alpha:W\to V$ such that
\[
        \|\widetilde{\varphi}_\alpha\|_{\pcb}
        \leq(\lambda+\eta)\|\varphi\|_{\pcb}
\]
and $\widetilde{\varphi}_\alpha|_{W_0}\to\varphi$ in the point-norm topology.
The net $\{P\widetilde{\varphi}_\alpha\}$ consists of maps from $W$ into $V_0$,
and $(P\widetilde{\varphi}_\alpha)|_{W_0}\to P\circ\varphi=\varphi$ in the
point-norm topology.  Moreover,
\[
        \|P\widetilde{\varphi}_\alpha\|_{\pcb}
        \leq \|P\|_{\pcb}\|\widetilde{\varphi}_\alpha\|_{\pcb}
        \leq \|P\|_{\pcb}(\lambda+\eta)\|\varphi\|_{\pcb}.
\]
If $P=0$, then $V_0=\{0\}$ and there is nothing to prove.  Otherwise, for arbitrary
$\varepsilon>0$, we take $\eta=\varepsilon/\|P\|_{\pcb}$. Then
\[
        \|P\widetilde{\varphi}_\alpha\|_{\pcb}
        \leq (\lambda\|P\|_{\pcb}+\varepsilon)\|\varphi\|_{\pcb},
\]
which is approximate $\lambda\|P\|_{\pcb}$-$p$-injectivity for $V_0$.
\end{proof}

A fundamental question is how these injectivity conditions behave under duality. The following theorem shows that for the dual $p$-operator space $V^*$, approximate $\lambda$-$p$-injectivity coincides with $\lambda$-$p$-injectivity.

\begin{theorem}\label{thm:dual-exact-approx}
Let $V$ be a $p$-operator space and let $\lambda\geq 1$.  Then $V^*$ is
approximately $\lambda$-$p$-injective if and only if $V^*$ is
$\lambda$-$p$-injective.
\end{theorem}

\begin{proof}
It is clear that $\lambda$-$p$-injectivity implies approximate $\lambda$-$p$-injectivity.

For the converse, let $W_0\subseteq W$ be $p$-operator spaces, and let $\varphi:W_0\to V^*$ be $p$-completely bounded.
If $\varphi=0$, the conclusion is obvious, so we assume $\|\varphi\|_{\pcb}>0$.

Let $\mathcal{D}$ be the set of triples $d=(\varepsilon,F,\delta)$, where
$0<\varepsilon\leq 1$, $F$ is a finite subset of $W_0$, and $\delta>0$.
We order $\mathcal{D}$ by declaring
\[
        (\varepsilon_1,F_1,\delta_1)\preceq(\varepsilon_2,F_2,\delta_2)
\]
if $\varepsilon_2\leq\varepsilon_1$, $F_1\subseteq F_2$, and
$\delta_2\leq\delta_1$.  By approximate $\lambda$-$p$-injectivity of $V^*$,
for each $d=(\varepsilon,F,\delta)$ there exists a $p$-completely bounded map
$\widetilde{\varphi}_d:W\to V^*$ such that
\[
        \|\widetilde{\varphi}_d\|_{\pcb}
        \leq(\lambda+\varepsilon)\|\varphi\|_{\pcb}
\]
and
\[
        \|\widetilde{\varphi}_d(x)-\varphi(x)\|<\delta,\qquad x\in F.
\]
For each $d\in\mathcal{D}$, we have
$\|\widetilde{\varphi}_d\|_{\pcb}\leq(\lambda+1)\|\varphi\|_{\pcb}$, so the net
lies in the closed ball
\[
        \mathcal{K}_M=\{S\in \CBp(W,V^*):\|S\|_{\pcb}\leq M\},
        \qquad M=(\lambda+1)\|\varphi\|_{\pcb}.
\]
By \cite[Proposition~4.9]{Daws2007}, there is a $p$-complete isometric isomorphism
\[
        \CBp(W,V^*) \cong (W \otimesp V)^*.
\]
Under this identification, the weak$^*$ and point-weak$^*$ topologies on
$\mathcal{K}_M$ coincide, since $W\otimes V$ is dense in $W\otimesp V$.
By the Banach--Alaoglu theorem, $\mathcal{K}_M$ is compact in the weak$^*$ topology,
and hence in the point-weak$^*$ topology.
We may choose a point-weak$^*$ limit $T:W\to V^*$ of a subnet
$\{\widetilde{\varphi}_{d_i}\}$.
For each $x\in W_0$, we have $\|\widetilde{\varphi}_{d_i}(x)-\varphi(x)\|\to 0$, so
$T(x)=\varphi(x)$ and $T|_{W_0}=\varphi$.

We now estimate $\|T\|_{\pcb}$. We take $m\in\mathbb{N}$ and $w=[w_{ij}]\in M_m(W)$ with $\|w\|\leq 1$.
Since $\widetilde{\varphi}_{d_i}\to T$ in the point-weak$^*$ topology, we have
$(\widetilde{\varphi}_{d_i})_m(w)\to T_m(w)$ in the weak$^*$ topology of $M_m(V^*)$.
By \cite[Section~5]{Daws2007}, for each $m\in\mathbb{N}$ there is a
$p$-complete isometric identification
\[
M_m(V^*)=\CBp(V,M_m)\cong (\mathcal{N}(\ell_p^m)\otimesp V)^*,
\]
where $\mathcal{N}(\ell_p^m)$ is the space of nuclear operators on $\ell_p^m$.
Hence $M_m(V^*)$ is a dual $p$-operator space, and the norm on $M_m(V^*)$
is weak$^*$ lower semicontinuous.
If $d_i=(\varepsilon_i,F_i,\delta_i)$, then
\[
        \|(\widetilde{\varphi}_{d_i})_m(w)\|
        \leq\|\widetilde{\varphi}_{d_i}\|_{\pcb}
        \leq(\lambda+\varepsilon_i)\|\varphi\|_{\pcb}.
\]
Passing to a subnet with $\varepsilon_i\to 0$, we obtain
\[
        \|T_m(w)\|\leq\liminf_i\|(\widetilde{\varphi}_{d_i})_m(w)\|
        \leq\liminf_i(\lambda+\varepsilon_i)\|\varphi\|_{\pcb}
        =\lambda\|\varphi\|_{\pcb}.
\]
We take the supremum over $m$ and over $w$ in the unit ball of $M_m(W)$
to obtain $\|T\|_{\pcb}\leq\lambda\|\varphi\|_{\pcb}$, so $V^*$ is
$\lambda$-$p$-injective.
\end{proof}

\begin{corollary}\label{cor:dual-p-2}
Let $V$ be an operator space ($p=2$) and let $\lambda\geq 1$.  Then
the following are equivalent:
\begin{enumerate}
    \item[\rm(1)] $V$ has the $\lambda$-ALLP;
    \item[\rm(2)] $V$ has the $\lambda$-LLP;
    \item[\rm(3)] $V^*$ is approximately $\lambda$-injective;
    \item[\rm(4)] $V^*$ is $\lambda$-injective.
\end{enumerate}
\end{corollary}

\begin{proof}
The equivalence (1)$\Leftrightarrow$(2) follows from
Theorem~\ref{thm:main} with $p=2$.
The equivalence (2)$\Leftrightarrow$(4) is established in
\cite{KyeRuan1999}.  The implication (2)$\Rightarrow$(3) is part of the chain in
\cite{Dong2007}.  Finally, (3)$\Leftrightarrow$(4) is the specialization of
Theorem~\ref{thm:dual-exact-approx} to $p=2$.
\end{proof}

\section{Dimension-dependent estimates}

\begin{proposition}\label{prop:fd-injective}
Let $V$ be a finite-dimensional $p$-operator space with $n=\dim V\geq 1$.
Then $V$ is $n$-$p$-injective.
\end{proposition}

\begin{proof}
We choose an Auerbach basis $v_1,\ldots,v_n$ of $V$ with coordinate functionals
$f_1,\ldots,f_n\in V^*$ satisfying $\|v_k\|=\|f_k\|=1$.  Let $W_0\subseteq W$ be $p$-operator spaces, and let $\psi:W_0\to V$ be $p$-completely bounded.
If $\psi=0$, we define $\widetilde\psi:W\to V$ to be the zero map.
We assume that $\psi\neq 0$, and we define $\psi_1=\psi/\|\psi\|_{\pcb}$, so $\|\psi_1\|_{\pcb}=1$.
By \cite[Lemma~4.2]{Daws2007}, each
$f_k\circ\psi_1:W_0\to\mathbb C$ satisfies
$\|f_k\circ\psi_1\|_{\pcb}=\|f_k\circ\psi_1\|\leq 1$.
Each $f_k\circ\psi_1$ extends to a functional $g_k:W\to\mathbb C$ with
$\|g_k\|\leq 1$ by the Hahn--Banach theorem, and
$\|g_k\|_{\pcb}=\|g_k\|$ by \cite[Lemma~4.2]{Daws2007}.  We define
$\widetilde\psi_1:W\to V$ by
\[
        \widetilde\psi_1(w)=\sum_{k=1}^n g_k(w)\,v_k,\qquad w\in W.
\]
For $w_0\in W_0$ and each $k$, we have $g_k(w_0)=f_k(\psi_1(w_0))$.  Since
$f_j(v_k)=\delta_{jk}$, we have $\psi_1(w_0)=\sum_k f_k(\psi_1(w_0))v_k$.
Hence $\widetilde\psi_1(w_0)=\sum_k g_k(w_0)v_k=\psi_1(w_0)$ for any
$w_0\in W_0$, and thus $\widetilde\psi_1|_{W_0}=\psi_1$.
We write $\widetilde\psi_1=\sum_{k=1}^n \theta_{v_k}\circ g_k$, where each map
$\theta_{v_k}:\mathbb{C}\to V$ is defined by $\theta_{v_k}(\alpha)=\alpha v_k$ and satisfies
$\|\theta_{v_k}\|_{\pcb}=\|v_k\|=1$. It follows that
\[
        \|\widetilde\psi_1\|_{\pcb}\leq
        \sum_{k=1}^n\|\theta_{v_k}\|_{\pcb}\|g_k\|_{\pcb}\leq n.
\]
Thus $\widetilde\psi=\|\psi\|_{\pcb}\,\widetilde\psi_1$ extends $\psi$ and satisfies
$\|\widetilde\psi\|_{\pcb}\leq n\|\psi\|_{\pcb}$.
\end{proof}

\begin{corollary}\label{cor:fd-equivalence}
Let $V$ be a finite-dimensional $p$-operator space with $n=\dim V\geq 1$.
Then $V$ has the $n$-$p$-LP, the $n$-$p$-LLP, and both $V$ and $V^*$ are
$n$-$p$-injective.
\end{corollary}

\begin{proof}
By Proposition~\ref{prop:fd-injective}, both $V$ and $V^*$ are $n$-$p$-injective.  Given $\varepsilon>0$, we choose $\eta>0$ with
$n(1+\eta)<n+\varepsilon$.  From Lemma~\ref{lem:correction} applied with $E=V$, for any
$p$-complete contraction $\varphi:V\to Y/W$ there exists $\widetilde\varphi:V\to Y$
such that $q\widetilde\varphi=\varphi$ and
$\|\widetilde\varphi\|_{\pcb}\leq n(1+\eta)<n+\varepsilon$.  Hence $V$ has the $n$-$p$-LP and therefore the $n$-$p$-LLP.
\end{proof}

\section{Injectivity of $B(\ell_p)$ and LLP of $\mathcal{N}(\ell_p)$}

We follow \cite{Daws2007} and write $\mathcal{N}(\ell_p)$ for the space of nuclear
operators on $\ell_p$.
By \cite[Lemma~5.1]{Daws2007}, $\mathcal{N}(\ell_p)'=B(\ell_p)$ $p$-completely isometrically.
For $n\in\mathbb{N}$, let $\iota_n:\ell_p^n\to\ell_p$ and $p_n:\ell_p\to\ell_p^n$ be the
coordinate inclusion and projection, and let $P_n:\ell_p\to\ell_p$ be the
corresponding multiplication projection onto the first $n$ coordinates.
As in the proof of \cite[Proposition~5.2]{Daws2007}, we define
\[
        j_n:\mathcal{N}(\ell_p^n)\to \mathcal{N}(\ell_p),\quad \tau\mapsto \iota_n\tau p_n,
        \qquad
        \mathcal{P}_n:\mathcal{N}(\ell_p)\to \mathcal{N}(\ell_p^n),\quad \sigma\mapsto p_n\sigma\iota_n.
\]
Then $j_n$ is a $p$-complete isometry, $\mathcal{P}_n$ is a $p$-complete quotient map,
$\mathcal{P}_n j_n=\mathrm{id}_{\mathcal{N}(\ell_p^n)}$, and $j_n\mathcal{P}_n$ is a completely contractive projection of $\mathcal{N}(\ell_p)$ onto
$j_n(\mathcal{N}(\ell_p^n))$.
For $\omega\in \mathcal{N}(\ell_p)$ we have $P_n\omega P_n=j_n(\mathcal{P}_n(\omega))$ as
operators on $\ell_p$.

\begin{lemma}\label{lem:Qn-approx-identity}
Let $1<p<\infty$. For each $n\in\mathbb{N}$, let $P_n:\ell_p\to\ell_p$ be the projection onto the first $n$ coordinates, and let $Q_n:B(\ell_p)\to B(\ell_p)$ be given by $Q_n(T)=P_nTP_n$.
Then $Q_n\to\mathrm{id}_{B(\ell_p)}$ in the point-weak$^*$ topology.
\end{lemma}

\begin{proof}
For $T\in B(\ell_p)$ and $\omega\in \mathcal{N}(\ell_p)$, we have
\[
        \langle Q_n(T),\omega\rangle=\langle T,P_n\omega P_n\rangle.
\]

To show that $Q_n\to\mathrm{id}_{B(\ell_p)}$ in the point-weak$^*$ topology,
it suffices to prove that $P_n\omega P_n\to\omega$ in norm for each $\omega\in \mathcal{N}(\ell_p)$.
Since $P_n$ is a contractive projection on $\ell_p$, the ideal property of
nuclear operators gives
\[
        \|P_n\tau P_n\|\leq\|\tau\|\qquad\text{for all }\tau\in \mathcal{N}(\ell_p).
\]
As in the proof of \cite[Proposition~5.2]{Daws2007}, the subspace
\[
        \bigcup_{m=1}^\infty j_m(\mathcal{N}(\ell_p^m))
\]
is norm dense in $\mathcal{N}(\ell_p)$.
Given $\omega\in \mathcal{N}(\ell_p)$ and $\varepsilon>0$, we choose
$\tau\in \bigcup_{m=1}^\infty j_m(\mathcal{N}(\ell_p^m))$ with
\[
        \|\omega-\tau\|<\varepsilon/2.
\]
There exists $N\in\mathbb{N}$ such that
\[
        P_n\tau P_n=\tau\qquad\text{for all }n\geq N.
\]
For $n\geq N$,
\[
        \|P_n\omega P_n-\omega\|
        \leq\|P_n(\omega-\tau)P_n\|+\|P_n\tau P_n-\tau\|+\|\tau-\omega\|
        \leq 2\|\omega-\tau\|<\varepsilon.
\]
Hence $P_n\omega P_n\to\omega$ in norm for each $\omega\in \mathcal{N}(\ell_p)$, so
$\langle Q_n(T),\omega\rangle\to\langle T,\omega\rangle$ for all $T\in B(\ell_p)$ and
$\omega\in \mathcal{N}(\ell_p)$, and hence $Q_n\to\mathrm{id}_{B(\ell_p)}$ in the point-weak$^*$ topology.
\end{proof}

\begin{theorem}\label{thm:B-lp-Mn-injective-equiv}
Let $1 < p < \infty$. The following are equivalent:
\begin{enumerate}
        \item[\rm(1)] $B(\ell_p)$ is $p$-injective.
        \item[\rm(2)] For every $n \geq 1$, $M_n=B(\ell_p^n)$ is $p$-injective.
\end{enumerate}
\end{theorem}

\begin{proof}
(1)$\Rightarrow$(2):
For each $n\geq 1$, we let $W_0\subseteq W$ be $p$-operator spaces and we let
$\theta:W_0\to M_n$ be $p$-completely bounded.
Let
\[
        \jmath_n:M_n\to B(\ell_p),\qquad \jmath_n(a)=\iota_n a p_n,
        \qquad
        R_n:B(\ell_p)\to M_n,\qquad R_n(T)=p_nT\iota_n .
\]
The maps $\jmath_n$ and $R_n$ are $p$-completely contractive and
$R_n\jmath_n=\mathrm{id}_{M_n}$.
By $p$-injectivity of $B(\ell_p)$, there exists $\Phi:W\to B(\ell_p)$ such that
$\Phi|_{W_0}=\jmath_n\circ\theta$ and
\[
        \|\Phi\|_{\pcb}\leq\|\theta\|_{\pcb}.
\]
For $w_0\in W_0$, $(R_n\circ\Phi)(w_0)=R_n(\jmath_n(\theta(w_0)))=\theta(w_0)$, so
$(R_n\circ\Phi)|_{W_0}=\theta$, and
\[
        \|R_n\circ\Phi\|_{\pcb}\leq\|\theta\|_{\pcb}.
\]
Thus $M_n$ is $p$-injective.

(2)$\Rightarrow$(1):
Let $W_0\subseteq W$ be $p$-operator spaces and let
$\psi:W_0\to B(\ell_p)$ be $p$-completely bounded.
For each $n\geq 1$, we let $P_n:\ell_p\to\ell_p$ be the projection onto the first
$n$ coordinates, and we define
\[
        Q_n:B(\ell_p)\to P_n B(\ell_p) P_n \cong M_n,\qquad Q_n(T)=P_nTP_n.
\]
Then $Q_n$ is the canonical weak$^*$ continuous $p$-completely contractive
finite-rank projection.
We define
\[
        \varphi_n=Q_n\circ\psi:W_0\to P_n B(\ell_p) P_n.
\]
By~(2), $P_n B(\ell_p) P_n$ is $p$-injective.
Thus, there exists a $p$-completely bounded map
$\widetilde\varphi_n:W\to P_n B(\ell_p) P_n$ with
\[
        \widetilde\varphi_n|_{W_0}=\varphi_n
        \qquad\text{and}\qquad
        \|\widetilde\varphi_n\|_{\pcb}\leq\|\varphi_n\|_{\pcb}\leq\|\psi\|_{\pcb}.
\]
We regard each $\widetilde\varphi_n$ as a $p$-completely bounded map $W\to B(\ell_p)$.
By \cite[Proposition~4.9]{Daws2007},
\[
        \CBp(W,B(\ell_p))\cong (W\otimesp \mathcal{N}(\ell_p))^*.
\]
Since $\|\widetilde\varphi_n\|_{\pcb}\leq\|\psi\|_{\pcb}$ for all $n$, the net is contained in the closed ball
\[
        \mathcal{B}_\psi=\{\Psi\in\CBp(W,B(\ell_p)):\|\Psi\|_{\pcb}\leq\|\psi\|_{\pcb}\}.
\]
Since $\CBp(W,B(\ell_p))$ is the $p$-operator space dual of $W\otimesp \mathcal{N}(\ell_p)$,
$\mathcal{B}_\psi$ is weak$^*$ compact.
Hence $\{\widetilde\varphi_n\}$ admits a subnet $\{\widetilde\varphi_{n_i}\}$ with point-weak$^*$ limit
$\widetilde\varphi\in\mathcal{B}_\psi$.
By the weak$^*$ lower semicontinuity of the norm,
\[
        \|\widetilde\varphi\|_{\pcb}\leq\liminf_{i\to\infty}\|\widetilde\varphi_{n_i}\|_{\pcb}\leq\|\psi\|_{\pcb}.
\]
For each $w_0\in W_0$, point-weak$^*$ convergence gives
\[
        \widetilde\varphi(w_0)=\lim_{i\to\infty}\widetilde\varphi_{n_i}(w_0),
\]
and $\widetilde\varphi_{n_i}|_{W_0}=Q_{n_i}\circ\psi$ implies
\[
        \lim_{i\to\infty}\widetilde\varphi_{n_i}(w_0)
        =\lim_{i\to\infty}Q_{n_i}\psi(w_0).
\]
By Lemma~\ref{lem:Qn-approx-identity}, $Q_n\to\mathrm{id}_{B(\ell_p)}$ in the point-weak$^*$ topology.

For $T=\psi(w_0)$,
\[
        \lim_{i\to\infty}Q_{n_i}\psi(w_0)=\psi(w_0),
\]
and therefore $\widetilde\varphi(w_0)=\psi(w_0)$ for all $w_0\in W_0$.
Thus $\widetilde\varphi|_{W_0}=\psi$, and $B(\ell_p)$ is $p$-injective.
\end{proof}

\begin{lemma}\label{lem:N-lpn-bidual-cb}
Let $1<p<\infty$ and $n\in\mathbb{N}$. Let $Y$ be a $p$-operator space on $L_p$ spaces.
Then there is a $p$-completely isometric identification
\[
        \CBp(\mathcal{N}(\ell_p^n),Y)^{**} \cong \CBp(\mathcal{N}(\ell_p^n),Y^{**}).
\]
\end{lemma}

\begin{proof}
Since $\mathcal{N}(\ell_p^n)^*\cong M_n$ $p$-completely isometrically, by
\cite[\S4.2]{Daws2007} we have a $p$-completely isometric identification
\[
        \CBp(\mathcal{N}(\ell_p^n),Y^{**})
        \cong \big(\mathcal{N}(\ell_p^n) \otimesp Y^*\big)^*
        \cong \CBp(Y^*, M_n)
        \cong M_n(Y^{**}).
\]
Because $Y$ is a $p$-operator space on $L_p$ spaces,
\[
        Y\subseteq Y^{**}
\]
$p$-completely isometrically by \cite[Proposition~4.4]{Daws2007}.
Thus
\[
        M_n(Y)\subseteq M_n(Y^{**}),
        \qquad
        \CBp(\mathcal{N}(\ell_p^n),Y)\subseteq \CBp(\mathcal{N}(\ell_p^n),Y^{**}),
\]
by restricting the above identification to these subspaces, we have
\[
        \CBp(\mathcal{N}(\ell_p^n),Y)\cong M_n(Y).
\]
Taking biduals and applying \cite[Theorem~3.6]{ZD14}, we obtain
\[
        \CBp(\mathcal{N}(\ell_p^n),Y)^{**}
        \cong M_n(Y)^{**}
        \cong M_n(Y^{**})
        \cong \CBp(\mathcal{N}(\ell_p^n),Y^{**}).
\]
\end{proof}

\begin{lemma}\label{lem:approx-from-bidual}
Let $E$ be a finite-dimensional $p$-operator space and $W\subseteq Y$ be $p$-operator spaces.
Let $q:Y\to Y/W$ denote the $p$-complete quotient map, and let $\varphi:E\to Y/W$ be a $p$-complete contraction.
Suppose that $\CBp(E,Y)^{**} \cong \CBp(E,Y^{**})$ $p$-completely isometrically, and that there exists a $p$-complete contraction $T:E\to Y^{**}$ such that for all $x\in E$ and $\mu\in (Y/W)^*$,
\[
        T(x)(\mu)=\mu(\varphi(x)).
\]
Then for any $\varepsilon>0$, there exists $\widetilde\psi\in\CBp(E,Y)$ such that $\|\widetilde\psi\|_{\pcb}\leq 1$ and $\|q\widetilde\psi-\varphi\|<\varepsilon$.
\end{lemma}

\begin{proof}
By Goldstine's theorem, the closed unit ball of $\CBp(E,Y)$ is
weak$^*$ dense in the closed unit ball of $\CBp(E,Y)^{**}$.
Thus there exists a net $\{\psi_\alpha\}\subseteq\CBp(E,Y)$ with
$\|\psi_\alpha\|_{\pcb}\leq 1$ such that
$\kappa_Y\circ\psi_\alpha\longrightarrow T$
in the point-weak$^*$ topology of $\CBp(E,Y^{**})$.

For $x\in E$ and $\mu\in (Y/W)^*$, we have
\[
        \mu\big(q\psi_\alpha(x)\big)
        =\mu(\psi_\alpha(x))
        =\kappa_Y(\psi_\alpha(x))(\mu)
        \longrightarrow T(x)(\mu)
        =\mu(\varphi(x)),
\]
so $q\psi_\alpha(x)\to\varphi(x)$ in the weak topology of $Y/W$ for each
$x\in E$.

Since $E$ is finite-dimensional, we choose a basis $x_1,\ldots,x_k$ with $k=\dim E$.
Let
\[
        C_E=\max\Bigl\{\sum_{j=1}^k|\alpha_j|:
        \Bigl\|\sum_{j=1}^k\alpha_j x_j\Bigr\|=1\Bigr\}<\infty.
\]
By Mazur's theorem, there exists a net of convex combinations
\[
        \widetilde\psi_\beta=\sum_\alpha t_\alpha^\beta \psi_\alpha,
\]
with $\|\widetilde\psi_\beta\|_{\pcb}\leq 1$, such that
\[
        \max_{1\leq j\leq k}
        \bigl\|q\widetilde\psi_\beta(x_j)-\varphi(x_j)\bigr\|
        \longrightarrow 0
\]
when $\beta\to\infty$.
For $y=\sum_{j=1}^k\alpha_j x_j\in E$ with $\|y\|=1$, we have
\begin{align*}
        \bigl\|q\widetilde\psi_\beta(y)-\varphi(y)\bigr\|
        &\leq\sum_{j=1}^k|\alpha_j|
        \max_{1\leq j\leq k}
        \bigl\|q\widetilde\psi_\beta(x_j)-\varphi(x_j)\bigr\| \\
        &\leq C_E
        \max_{1\leq j\leq k}
        \bigl\|q\widetilde\psi_\beta(x_j)-\varphi(x_j)\bigr\|.
\end{align*}
Hence $\|q\widetilde\psi_\beta-\varphi\|\to 0$ when $\beta\to\infty$.
Thus, for any given $\varepsilon>0$, we have
\[
        \|q\widetilde\psi_\beta-\varphi\|<\varepsilon
\]
whenever $\beta$ is sufficiently large.
\end{proof}


\begin{proposition}\label{prop:N-lpn-ALLP}
Let $1<p<\infty$ and $n\in\mathbb{N}$.
If $M_n$ is $p$-injective, then $\mathcal{N}(\ell_p^n)$ has the $p$-ALLP for $p$-operator spaces on $L_p$ spaces.
\end{proposition}

\begin{proof}
We take $p$-operator spaces $W\subseteq Y$ on $L_p$ spaces,
let $q:Y\to Y/W$ denote the $p$-complete quotient map, and let
$\varphi:\mathcal{N}(\ell_p^n)\to Y/W$ be a $p$-complete contraction.
Let $\varepsilon>0$.

By \cite[Lemma~4.5]{Daws2007}, the adjoint map
\[
        \varphi^*:(Y/W)^*\longrightarrow \mathcal{N}(\ell_p^n)^*
\]
is a $p$-complete contraction.
Under the annihilator identification $(Y/W)^*\cong W^\perp\subseteq Y^*$
(\cite[Lemma~4.6]{Daws2007}), where
\[
        W^\perp=\{\mu\in Y^*:\mu|_W=0\},
\]
we view $\varphi^*$ as a $p$-complete contraction $W^\perp\to \mathcal{N}(\ell_p^n)^*$.
Since
\[
        \mathcal{N}(\ell_p^n)^*\cong B(\ell_p^n)=M_n,
\]
the assumed $p$-injectivity of $M_n$ implies that there exists a $p$-completely
contractive extension
\[
        \Phi_n:Y^*\longrightarrow M_n
\]
of $\varphi^*$.

By \cite[Lemma~4.5]{Daws2007}, its adjoint $\Phi_n^*:M_n^*\to Y^{**}$ is a
$p$-complete contraction.
Let $\kappa:\mathcal{N}(\ell_p^n)\to \mathcal{N}(\ell_p^n)^{**}=M_n^*$ be the canonical $p$-complete isometry.
We define
\[
        T_n = \Phi_n^*\circ\kappa \colon \mathcal{N}(\ell_p^n) \longrightarrow Y^{**},
\]
a $p$-complete contraction.
For $v\in \mathcal{N}(\ell_p^n)$ and $\mu\in W^\perp$, we have
\begin{equation}\label{eq:T-matches-phi-lem}
        T_n(v)(\mu)
        =\kappa(v)\big(\Phi_n(\mu)\big)
        =\Phi_n(\mu)(v)
        =\varphi^*(\mu)(v)
        =\mu\big(\varphi(v)\big).
\end{equation}
By Lemma~\ref{lem:N-lpn-bidual-cb}, we have
\[
        \CBp(\mathcal{N}(\ell_p^n),Y)^{**} \cong \CBp(\mathcal{N}(\ell_p^n),Y^{**})
\]
$p$-completely isometrically.
Together with the identity above and the identification $(Y/W)^*\cong W^\perp$, the map $T_n$ satisfies the hypotheses of Lemma~\ref{lem:approx-from-bidual}.
Hence there exists $\widetilde\psi\in\CBp(\mathcal{N}(\ell_p^n),Y)$ such that $\|\widetilde\psi\|_{\pcb}\leq 1$ and $\|q\widetilde\psi-\varphi\|<\varepsilon$.
Since $\varepsilon>0$ is arbitrary, this verifies Definition~\ref{def:lambda-allp} with $\lambda=1$.
Thus $\mathcal{N}(\ell_p^n)$ has the $p$-ALLP for $p$-operator spaces on $L_p$ spaces.
\end{proof}

\begin{lemma}\label{lem:Pn-approx-identity}
Let $1<p<\infty$.
For each $n\in\mathbb{N}$, let $\iota_n:\ell_p^n\to\ell_p$ and $p_n:\ell_p\to\ell_p^n$ be the
coordinate inclusion and projection, and let $P_n:\ell_p\to\ell_p$ be the projection onto the
first $n$ coordinates.
We define
\[
        j_n:\mathcal{N}(\ell_p^n)\to \mathcal{N}(\ell_p),\quad \tau\mapsto \iota_n\tau p_n,
        \qquad
        \mathcal{P}_n:\mathcal{N}(\ell_p)\to \mathcal{N}(\ell_p^n),\quad \sigma\mapsto p_n\sigma\iota_n.
\]
Then $j_n\mathcal{P}_n\to \mathrm{id}_{\mathcal{N}(\ell_p)}$ in norm.
In particular, if $E\subseteq \mathcal{N}(\ell_p)$ is a finite-dimensional subspace, then $j_n\mathcal{P}_n|_E\to\mathrm{id}_E$ in norm.
\end{lemma}

\begin{proof}
For $\omega\in \mathcal{N}(\ell_p)$ we have
$j_n\mathcal{P}_n(\omega)=P_n\omega P_n$ as operators on $\ell_p$.
Since $P_n$ is a contractive projection on $\ell_p$, the ideal property of
nuclear operators gives $\|P_n\tau P_n\|\leq\|\tau\|$ for all
$\tau\in \mathcal{N}(\ell_p)$.
As in the proof of \cite[Proposition~5.2]{Daws2007}, the subspace $\bigcup_{m=1}^\infty j_m(\mathcal{N}(\ell_p^m))$ is
norm dense in $\mathcal{N}(\ell_p)$.
Given $x\in \mathcal{N}(\ell_p)$ and $\delta>0$, we choose
$\tau\in \bigcup_{m=1}^\infty j_m(\mathcal{N}(\ell_p^m))$ with $\|x-\tau\|<\delta/2$, and then $M\in\mathbb{N}$ with $P_n\tau P_n=\tau$ for
all $n\geq M$.
For $n\geq M$,
\[
        \|j_n\mathcal{P}_n(x)-x\|
        =\|P_nxP_n-x\|
        \leq\|P_n(x-\tau)P_n\|+\|P_n\tau P_n-\tau\|+\|\tau-x\|
        \leq 2\|x-\tau\|<\delta.
\]
Hence $j_n\mathcal{P}_n\to \mathrm{id}_{\mathcal{N}(\ell_p)}$ in norm.
Since $E$ is finite-dimensional, $j_n\mathcal{P}_n|_E\to\mathrm{id}_E$ in norm.
\end{proof}

\begin{proposition}\label{prop:inj-to-llp}
Let $1<p<\infty$.  Suppose that $B(\ell_p)$ is $p$-injective.
Then $\mathcal{N}(\ell_p)$ has the $p$-LLP for $p$-operator spaces on $L_p$ spaces.
\end{proposition}

\begin{proof}
We have
\[
        \mathcal{N}(\ell_p)^*=B(\ell_p)
\]
$p$-completely isometrically by \cite[Lemma~5.1]{Daws2007}.
By Remark~\ref{rem:main-on-Lp}, it suffices to show that $\mathcal{N}(\ell_p)$ has the
$p$-ALLP for $p$-operator spaces on $L_p$ spaces.

We take $p$-operator spaces $W\subseteq Y$ on $L_p$ spaces,
let $q:Y\to Y/W$ denote the $p$-complete quotient map, let
$\varphi:\mathcal{N}(\ell_p)\to Y/W$ be a $p$-complete contraction, let
$E\subseteq \mathcal{N}(\ell_p)$ be finite-dimensional, and let $\varepsilon>0$.

As above, $\mathcal{P}_n$ is a $p$-complete quotient map (since it has a $p$-completely isometric right inverse $j_n$).
By Lemma~\ref{lem:Pn-approx-identity}, $j_n\mathcal{P}_n|_E\to\mathrm{id}_E$ in norm.

We may assume that $E\neq\{0\}$ (otherwise the condition is trivial).
We choose $n$ such that
\[
        \|j_n\mathcal{P}_n|_E-\mathrm{id}_E\|<\min\{1,\varepsilon/2\}.
\]
In particular, since $\|\mathrm{id}_E\|=1$, this guarantees that $\|\mathcal{P}_n|_E\|\neq 0$.
We define
\[
        \varphi_n=\varphi\circ j_n:\mathcal{N}(\ell_p^n)\to Y/W .
\]
Then $\|\varphi_n\|_{\pcb}\leq 1$ and
\[
        \|\varphi_n\mathcal{P}_n|_E-\varphi|_E\|
        =\|\varphi j_n\mathcal{P}_n|_E-\varphi|_E\|<\varepsilon/2 .
\]

Since $B(\ell_p)$ is $p$-injective, $M_n$ is $p$-injective by Theorem~\ref{thm:B-lp-Mn-injective-equiv}.
By Proposition~\ref{prop:N-lpn-ALLP}, $\mathcal{N}(\ell_p^n)$ has the $p$-ALLP for $p$-operator spaces on $L_p$ spaces.
Applying this to the $p$-complete contraction $\varphi_n:\mathcal{N}(\ell_p^n)\to Y/W$ and the finite-dimensional subspace $\mathcal{P}_n(E)\subseteq \mathcal{N}(\ell_p^n)$, there exists a $p$-completely bounded map
$\widetilde\psi:\mathcal{N}(\ell_p^n)\to Y$ such that $\|\widetilde\psi\|_{\pcb}\leq 1$ and
\[
        \|q\widetilde\psi|_{\mathcal{P}_n(E)}-\varphi_n|_{\mathcal{P}_n(E)}\|<\frac{\varepsilon}{2\|\mathcal{P}_n|_E\|}.
\]
We define
\[
        \widehat\psi=\widetilde\psi\circ \mathcal{P}_n|_E:E\to Y.
\]
Then $\|\widehat\psi\|_{\pcb}\leq 1$, and we have
\begin{align*}
        \|q\widehat\psi|_E-\varphi|_E\|
        &\leq
        \|q\widetilde\psi \mathcal{P}_n|_E-\varphi_n\mathcal{P}_n|_E\|
        +\|\varphi_n\mathcal{P}_n|_E-\varphi|_E\| \\
        &\leq
        \|q\widetilde\psi|_{\mathcal{P}_n(E)}-\varphi_n|_{\mathcal{P}_n(E)}\| \cdot \|\mathcal{P}_n|_E\|
        +\|\varphi_n\mathcal{P}_n|_E-\varphi|_E\| \\
        &< \varepsilon/2 + \varepsilon/2 = \varepsilon.
\end{align*}
Since $\varepsilon>0$ is arbitrary, this verifies Definition~\ref{def:lambda-allp} with $\lambda=1$.
Thus $\mathcal{N}(\ell_p)$ has the $p$-ALLP for $p$-operator spaces on $L_p$ spaces, and
Remark~\ref{rem:main-on-Lp} then yields the $p$-LLP in the same sense.
\end{proof}

\begin{thebibliography}{99}

\bibitem{AnLeeRuan2010}
G. An, J.-J. Lee, and Z.-J. Ruan,
\emph{On $p$-approximation properties for $p$-operator spaces},
J. Funct. Anal. \textbf{259} (2010), 933--974.

\bibitem{Ble92}
D. P. Blecher,
\emph{The standard dual of an operator space},
Pacific J. Math. \textbf{153} (1992), 15--30.

\bibitem{CE76}
M.-D. Choi and E. G. Effros,
\emph{Separable nuclear $C^*$-algebras and injectivity},
Duke Math. J. \textbf{43} (1976), 309--322.

\bibitem{CE77b}
M.-D. Choi and E. G. Effros,
\emph{Nuclear $C^*$-algebras and injectivity: the general case},
Indiana Univ. Math. J. \textbf{44} (1977), 443--446.

\bibitem{Co76}
A. Connes,
\emph{Classification of injective factors. Cases $II_1$, $II_\infty$, $III_\lambda$, $\lambda \ne 1$},
Ann. of Math. (2) \textbf{104} (1976), 73--115.

\bibitem{Daws2007}
M. Daws,
\emph{$p$-Operator spaces and Fig\`a--Talamanca--Herz algebras},
J. Operator Theory \textbf{63} (2010), 47--83.

\bibitem{Dong2007}
Z. Dong,
\emph{On the approximating local lifting property for operator spaces},
J. Math. Anal. Appl. \textbf{336} (2007), 1054--1060.

\bibitem{Dong2009}
Z. Dong,
\emph{On the local lifting properties of operator spaces},
Canad. J. Math. \textbf{61} (2009), 1262--1278.

\bibitem{DR07}
Z. Dong and Z.-J. Ruan,
\emph{Weak$^*$ exactness for dual operator spaces},
J. Funct. Anal. \textbf{253} (2007), 373--397.

\bibitem{EOR01}
E. G. Effros, N. Ozawa, and Z.-J. Ruan,
\emph{On injectivity and nuclearity for operator spaces},
Duke Math. J. \textbf{110} (2001), 489--521.

\bibitem{ER00}
E. G. Effros and Z.-J. Ruan,
\emph{Operator Spaces},
London Mathematical Society Monographs. New Series, vol. 23, The Clarendon Press, Oxford University Press, New York, 2000.

\bibitem{Gr55}
A. Grothendieck,
\emph{Produits tensoriels topologiques et espaces nucl\'eaires},
Mem. Amer. Math. Soc. \textbf{16} (1955).

\bibitem{KyeRuan1999}
S.-H. Kye and Z.-J. Ruan,
\emph{On the local lifting property for operator spaces},
J. Funct. Anal. \textbf{168} (1999), 355--379.

\bibitem{Lee2014}
J.-J. Lee,
\emph{Conditions $C_p$, $C'_p$, and $C''_p$ for $p$-operator spaces},
Oper. Matrices \textbf{8} (2014), 1079--1093.

\bibitem{Lee2015}
J.-J. Lee,
\emph{Hahn--Banach type extension theorems on $p$-operator spaces},
Oper. Matrices \textbf{9} (2015), 675--685.

\bibitem{LeMerdy1996}
C. Le Merdy,
\emph{Factorization of $p$-completely bounded multilinear maps},
Pacific J. Math. \textbf{172} (1996), 187--213.

\bibitem{Pisier1990}
G. Pisier,
\emph{Completely bounded maps between sets of Banach space operators},
Indiana Univ. Math. J. \textbf{39} (1990), 249--277.

\bibitem{Ru89}
Z.-J. Ruan,
\emph{Subspaces of $C^*$-algebras},
J. Funct. Anal. \textbf{76} (1988), 217--230.

\bibitem{Ryan2002}
B. Ryan,
\emph{Introduction to Tensor Products of Banach Spaces},
Springer, London, 2002.

\bibitem{ZD14}
Y. Zhao and Z. Dong,
\emph{$p$-mapping spaces for $p$-operator spaces on $L_p$ spaces},
Illinois J. Math. \textbf{58} (2014), 785--805.

\end{thebibliography}

\noindent
\begin{minipage}{\textwidth}
\noindent\textbf{Yafei Zhao}\\
Department of Mathematics\\
Zhejiang International Studies University\\
Hangzhou 310023, China\\
E-mail address: \texttt{yfzhao@zisu.edu.cn}

\medskip
\noindent\textbf{Zhe Dong}\\
School of Mathematical Sciences\\
Zhejiang University\\
Hangzhou 310058, China\\
E-mail address: \texttt{dongzhe@zju.edu.cn}
\end{minipage}

\end{document}
